\scanpage{108}

$K$ extension tot. réelle [i.e. si $K$ défini par $H\subset\pi_1(\mathbb Q)$, $H$
ne contient aucun conjugué de $\tau$,
$\tau$ la conjugaison ordinaire]

A) $L$ extension quadratique, $L=K(\unsure{e|x})$, $x^2+ax+b=0$ \quad$(a,b\in K)$.
\inserted{$\sigma$ l'élt non-trivial de $\mathrm{Gal}(L/K)$, alors $\sigma x=\bar x$ pour tt $x$}
\margin{$a=x+\deleted{\unsure{x}}\bar x$, $b=x\deleted{\unsure{x}}\bar x$}
Conditions équivalentes :

\diagram{\begin{tikzcd}[ampersand replacement=\&,row sep=small]
{\text{$L$ tot. imaginaire}}\arrow[d,Leftrightarrow]\\
{\begin{gathered}\text{Pour tt plongmt de $K$ dans $\mathbb R$, }\Delta=\deleted{\unsure{b}}a^2-4b\\
=(x-\deleted{\unsure{x}}\bar x)^2\text{ est }<\deleted{\unsure{\geq}}0,\ \text{i.e. $\Delta$ est}\\
\text{tot. \deleted{positif} \inserted{str.} négatif}\end{gathered}}\arrow[d,Leftrightarrow]\\
{\deleted{\begin{gathered}\text{$L$ est stable par $\unsure{x|t}\mapsto\bar x$, et}\\
(x-\bar x)^2\text{ est tot. positif}\end{gathered}}}
\end{tikzcd}}

\deleted{Alors $\sigma$ est la conj. complexe \unsure{ord.}}

B) Soit $x\in\overline{\mathbb Q}\subset\mathbb C$. Conditions équivalentes

\diagram{\begin{tikzcd}[ampersand replacement=\&,row sep=small]
{1)\ \left\{\begin{aligned}&x+\bar x\ [\text{et }x\bar x\text{ sont}]\ \inserted{\text{conséquence des autres, car }x\bar x=\frac14[(x+\bar x)^2-(x-\bar x)^2]}\ \text{tot. réels}\\
&(x-\bar x)^2\text{ est tot. \deleted{\unsure{positif $>0$|$\geq0$} positif \unsure{$<0$}} \inserted{str.} négatif}\end{aligned}\right.}\arrow[d,Leftrightarrow]\\
{2)\quad\begin{gathered}\deleted{\text{$x$ est quadr.}}\ \exists\text{ un $K$ tot. réel et une}\\
\text{extension $L$ tot. imaginaire de $K$, dont}\\
\text{$x$ soit un générateur.}\end{gathered}}\arrow[d,Leftrightarrow]\\
{2')\quad\deleted{\unsure{\exists}}\,x\in\mathbb Q'(i)-\mathbb Q',\quad\text{i.e.}}
\end{tikzcd}}

En effet, $2)\Rightarrow1)$ en vertu de A), et pour $1)\Rightarrow2)$, il suffit
de prendre $K=\mathbb Q(a,b)$ où $a=x+\bar x$, $b=x\bar x$, et appliquer
encore A).

\uline{Remarque} Cel\unsure{à|a} prouve que les
conditions de 1) sont stables par conjugaison, ce qui
n'était pas évident autrement.

\deleted{3) $x\bar x$ est tot. str. positif
et de même pour $x_\alpha\bar x_\alpha$ pour tt conjugué de $x$}

C) Soit $x\in\overline{\mathbb Q}\subset\mathbb C$. Conditions équivalentes

\diagram{\begin{tikzcd}[ampersand replacement=\&,row sep=small]
{1)\quad x+\bar x\text{ tot. réel},\quad(x-\bar x)^2\text{ tot. négatif}}\arrow[d,Leftrightarrow]\\
{2)\quad\exists\ K,L\text{ comme plus haut},\quad x\in\deleted{K}L}\arrow[d,Leftrightarrow]\\
{3)\quad\unsure{a|x}\in\mathbb Q'(i).}\arrow[d,Rightarrow,bend right=28]\\
{4)\quad x\bar x\text{ tot. positif, \deleted{et de même} \deleted{\unsure{équiv.|égal}} \inserted{ainsi que} }x_\alpha\bar x_\alpha\text{ pour tt conjugué $x_\alpha$ de $x$.}}\arrow[u,Rightarrow,"{?}"',bend left=55]
\end{tikzcd}}

% archive sheet number: 107
